\section{Continuous equivalent layers on a plane}
\label{eq:planar_eqlayer}

An equivalent layer is loosely defined as a physical property distribution that satisfactorily 
reproduces potential-field data and enables potential-field transformations. 
In the literature, equivalent layers are typically assumed to lie in source-free regions where 
Laplace’s equation holds. Consequently, it is presumed that the equivalent layer does not intersect 
the true sources. However, in practical applications, the positions of the true sources in the subsurface 
are often unknown, making this condition difficult to guarantee.
Empirical evidence suggests that equivalent-layer methods can still perform effectively 
when the layer intersects the sources \cite<e.g.,>[]{gonzalez_etal2022}.

In this section, we formally define \textit{continuous equivalent layers} 
on a horizontal surface $\mathcal{S}''$ that may either lie above or intersect the sources (Figure \ref{fig:Figure1}). 
Our goal is to leverage the properties of continuous layers to gain deeper insights into their discrete counterparts, 
which are explored in the companion paper 
``Planar equivalent layers for potential fields: 2. Geophysical implications and applications".

\begin{definition}[\textbf{Continuous equivalent layer}]
	Let us assume to have an observed potential field $d^{o}$ defined above 
	a horizontal surface $\mathcal{S}''$ with constant vertical coordinate $z''_{0}$. 
	Besides, consider (i) a horizontal surface $\mathcal{S}^{\dagger}$ that is located above 
	$\mathcal{S}''$, at a constant vertical coordinate $z^{\dagger}_{0}$, 
	and does not cross the sources and (ii) a harmonic kernel 
	$a$ that depends only on $\mathbf{r}$ and $\mathbf{r}'' \in \mathcal{S}''$, as well as 
	vanishes as the distance $\| \mathbf{r} - \mathbf{r}'' \|$ increases. 
	Then, for a chosen accuracy $\varepsilon > 0$, we may call $p_{\varepsilon}$ 
	a \underline{continuous equivalent layer} if its predicted potential field 
	\begin{equation}
		d_{\varepsilon}(\mathbf{r})
		= \int\limits_{-\infty}^{\infty}\int\limits_{-\infty}^{\infty}
		p_{\varepsilon}(\mathbf{r}'') \, a(\mathbf{r}, \mathbf{r}'') \:\: d\mathcal{S}''
		\: , \quad z < z''_{0} \: ,
		\label{eq:predicted-data-continuous-layer}
	\end{equation}
	satisfies 
	\begin{equation}
		\int\limits_{-\infty}^{\infty} \int\limits_{-\infty}^{\infty}
		\vert d^{o}(\mathbf{r}^{\dagger}) - 
		d_{\varepsilon}(\mathbf{r}^{\dagger}) \vert^{2}
		\:\: d\mathcal{S}^{\dagger} \le \varepsilon \: ,
		\quad z^{\dagger}_{0} < z^{t}_{0} \: , \quad z^{\dagger}_{0} < z''_{0} \: ,
		\label{eq:epsilon-bound-data}
	\end{equation}
	where $\mathbf{r} = (x, y, z)^{\top}$, 
	$\mathbf{r}^{\dagger} = (x^{\dagger}, y^{\dagger}, z^{\dagger}_{0})^{\top} \in \mathcal{S}^{\dagger}$ 
	and $\mathbf{r}'' = (x'', y'', z''_{0})^{\top} \in \mathcal{S}''$.
	\label{def:continuous-eqlayer}
\end{definition}

In Definition \ref{def:continuous-eqlayer}, $p_{\varepsilon}$ represents a 
continuous physical-property distribution that extends to infinity along the horizontal 
directions. 
Moreover, the kernel function $a$ is usually defined by $D_{a} \Psi$ 
as a generic partial derivative of $\Psi$ (equation \ref{eq:Psi}),
where $D_{a}$ is set in terms of the space-domain differential operator 
$D$ (equation \ref{eq:diff-operators-generic}).
It is important to stress here that the horizontal surface $\mathcal{S}''$ in 
Definition \ref{def:continuous-eqlayer} may cross the sources, according to 
Figure \ref{fig:Figure1}.

A particular continuous equivalent layer 
$p_{\varepsilon}$ leads to a transformed potential field given by:
\begin{equation}
	f_{\varepsilon}(\mathbf{r}^{\ddagger})
	= \int\limits_{-\infty}^{\infty}\int\limits_{-\infty}^{\infty}
	p_{\varepsilon}(\mathbf{r}'') \, t(\mathbf{r}^{\ddagger}, \mathbf{r}'') \:\: d\mathcal{S}'' \: , 
	\quad \mathbf{r}^{\ddagger} \in \mathcal{S}^{\ddagger} \: ,
	\quad \mathbf{r}'' \in \mathcal{S}'' \: ,
	\label{eq:transformed-field-continuous-layer}
\end{equation}
where $\mathbf{r}^{\ddagger}$ are points defined on a planar surface $\mathcal{S}^{\ddagger}$ located above 
$\mathcal{S}''$ (Figure \ref{fig:Figure1}); 
$f_{\varepsilon}$ is a transformed field that approximates $f$ (equations \ref{eq:transformed-field-space}); and
$t$ is the transformation kernel, which is usually defined by $D_{t} \Psi$ as a generic partial 
derivative of the potential function $\Psi$ (equation \ref{eq:Psi}), with $D_{t}$ given in terms of 
$D$ (equation \ref{eq:diff-operators-generic}).
In equation \ref{eq:transformed-field-continuous-layer}, we define $\mathbf{r}^{\ddagger}$ on 
a planar surface $\mathcal{S}^{\ddagger}$ located above $\mathcal{S}''$ (Figure \ref{fig:Figure1}), but this is not a 
restriction. They can coincide with the observation points or be arbitrarily defined above $\mathcal{S}''$.
